\sechead{3}{Variation of $G_\ell$ and $\widetilde{G}_\ell$ with $\ell$}%
  {VARIATION OF $G_\ell$ AND $\widetilde{G}_\ell$ WITH $\ell$}

\subsechead{3.1}{Preliminaries}

We keep the notations of the preceding paragraphs. For each prime number $\ell$, we denote by $\rho_\ell$ the homomorphism
\[
  \operatorname{Gal}(\overline{K}/K) \longrightarrow \operatorname{Aut}(T_\ell) \simeq \GL(2,\Z_\ell)
\]
defined by the action of $\operatorname{Gal}(\overline{K}/K)$ on $T_\ell$. The $\rho_\ell$'s define a homomorphism
\[
  \rho : \operatorname{Gal}(\overline{K}/K) \longrightarrow \prod_{\ell} \operatorname{Aut}(T_\ell),
\]
where the product is taken over the set of all prime numbers.

Let $G = \operatorname{Im}(\rho) \subsetneq \prod_{\ell} \operatorname{Aut}(T_\ell)$ and $G_\ell = \operatorname{Im}(\rho_\ell) \subsetneq \operatorname{Aut}(T_\ell)$, so that $G_\ell$ is the image of $G$ under the $\ell^{\text{th}}$ projection map. Let $\widetilde{G}_\ell$ be the image of $G_\ell$ in $\operatorname{Aut}(E_\ell) = \operatorname{Aut}(T_\ell/\ell T_\ell) \simeq \GL(2,\F_\ell)$.

\begin{lemma}
(1) The image of $G$ by $\det : \prod \operatorname{Aut}(T_\ell) \longrightarrow \prod \Z_\ell^{*}$ is open.
(2) For almost all $\ell$, $\det(G_\ell) = \Z_\ell^{*}$ and $\det(\widetilde{G}_\ell) = \F_\ell^{*}$.
\end{lemma}

We know (cf. chap.~\ref{chap:I}, \ref{sec:I.1.2}) that $\det(\rho_\ell) : \operatorname{Gal}(\overline{K}/K) \longrightarrow \Z_\ell^{*}$ is the character $\chi_\ell$ giving the action of $\operatorname{Gal}(\overline{K}/K)$ on $\ell^n$-th roots of unity. Hence $\det(G) \subsetneq \prod \Z_\ell^{*}$ is the Galois group $\operatorname{Gal}(K_c/K)$, where $K_c K$ is the extension of $K$ generated by all roots of unity. Since one knows that $\operatorname{Gal}(Q_c/Q) = \prod \Z_\ell^{*}$ (cf. for instance \cite{lang-algnum}, chap. IV) it follows that $\det(G)$ is the open subgroup of $\prod \Z_\ell^{*}$ corresponding to the field $K \cap Q_c$, hence (1). Assertion (2) follows from (1) and the definition of the product topology.

Assume now that $E$ has no complex multiplication. We know (cf. \ref{sec:IV.2.2}) that each $G_\ell$ is open in $\operatorname{Aut}(T_\ell)$. This does not \emph{a priori} imply that $G$ itself is open. However:

\begin{proposition}
The following properties are equivalent:
\begin{enumerate}
\item[$\mathrm{(i)}$] $G$ is open in $\prod_{\ell}\operatorname{Aut}(T_\ell)$.
\item[$\mathrm{(ii)}$] $G_\ell = \operatorname{Aut}(T_\ell)$ for almost all $\ell$.
\item[$\mathrm{(iii)}$] $\widetilde{G}_\ell = \operatorname{Aut}(E_\ell)$ for almost all $\ell$.
\item[$\mathrm{(iv)}$] $\widetilde{G}_\ell$ contains $\SL(E_\ell)$ for almost all $\ell$.
\end{enumerate}
\end{proposition}

The implications (i) $\Rightarrow$ (ii) $\Rightarrow$ (iii) $\Rightarrow$ (iv) are trivial. Implication (iv) $\Rightarrow$ (i) follows from the following group-theoretical result, whose proof will be given in section \ref{sec:IV.3.4} below:

\begin{stmt}{MAIN LEMMA}{}
Let $G$ be a closed subgroup of $\prod \GL(2,\Z_\ell)$ and let $G_\ell$ and $\widetilde{G}_\ell$ denote its images in $\GL(2,\Z_\ell)$ and $\GL(2,\F_\ell)$ as above. Assume:
\begin{enumerate}
\item[$\mathrm{(a)}$] $G_\ell$ is open in $\GL(2,\Z_\ell)$ for all $\ell$.
\item[$\mathrm{(b)}$] The image of $G$ by $\det : \prod \GL(2,\Z_\ell) \to \prod \Z_\ell^{*}$ is open.
\item[$\mathrm{(c)}$] $G_\ell$ contains $\SL(2,\F_\ell)$ for almost all $\ell$.
\end{enumerate}
Then $G$ is open in $\prod \GL(2,\Z_\ell)$.
\end{stmt}

\begin{remarks}
For each integer $n \geq 1$, let $E_n$ be the group of points of $E(\overline{K})$ of order dividing $n$, and let $\widetilde{G}_n$ be the image of the canonical map $\operatorname{Gal}(\overline{K}/K) \to \operatorname{Aut}(E_n) \simeq \GL(2,\Z/n\Z)$. One sees easily that property (i) above is equivalent to
\begin{enumerate}
\item[$\mathrm{(i')}$] The index of $\widetilde{G}_n$ in $\operatorname{Aut}(E_n)$ is bounded.
\end{enumerate}
\end{remarks}

\subsechead{3.2}{The case of a non-integral $j$}

\begin{theorem}
Assume that the modular invariant $j$ of $E$ is not an integer of $K$. Then $E$ enjoys the equivalent properties (i), (ii), (iii), (iv) of \ref{sec:IV.3.1}.
\end{theorem}

Since $j$ is not integral, we can choose a place $v$ of $K$ such that $v(j) < 0$. Let $q$ be the element of the local field $K_v$ which corresponds to $j$ by Tate's theory (cf. Appendix, A.1.1) and let $E_q$ be the corresponding elliptic curve over $K_v$. There is a finite extension $K'$ of $K_v$ over which $E$ and $E_q$ are isomorphic; one can even take for $K'$ either $K_v$ or a quadratic extension of $K_v$. Let $v'$ be the valuation of $K'$ which extends $v$; assume $v'$ is normalized so that $v'(K') = \Z$, and let $n = v'(q) = -v'(j)$. We have $n \geq 1$.

\begin{lemma}[1]\label{lem:IV.3.2.1}
Let $\ell$ not divide $n$, and let $I_{v,\ell}$ be the inertia subgroup of $\widetilde{G}_\ell$ corresponding to some extension of $v$ to $\overline{K}$. Then $I_{v,\ell}$ contains a transvection, i.e.\ an element whose matrix form is $\begin{pmatrix} 1 & 1 \\ 0 & 1 \end{pmatrix}$ for a suitable $\F_\ell$-basis of $E_\ell$.
\end{lemma}

This is true for the curve $E_q$ over $K'$, cf. Appendix, A.1.5. The result for $E$ follows from the isomorphism $E_{/K'} \simeq E_{q/K'}$.

\begin{lemma}[2]\label{lem:IV.3.2.2}
Let $H$ be a subgroup of $\GL(2,\F_\ell)$ which acts irreducibly on $\F_\ell \times \F_\ell$ and which contains a transvection. Then $H$ contains $\SL(2,\F_\ell)$.
\end{lemma}

\begin{proof}[Proof of Lemma~\ref{lem:IV.3.2.2}]
For any transvection $s \in H$, let $D_s$ be the unique one-dimensional subspace of $\F_\ell \times \F_\ell$ which is fixed by $s$. If all such lines were the same, the line so defined would be stable by $H$, and $H$ would not be irreducible. Hence there are transvections $s,s' \in H$ such that $D_s \neq D_{s'}$. If we choose a suitable basis $(e,e')$ of $\F_\ell \times \F_\ell$, this means that the matrix forms of $s,s'$ are
\[
  s = \begin{pmatrix} 1 & 1 \\ 0 & 1 \end{pmatrix}, \quad
  s' = \begin{pmatrix} 1 & 0 \\ 1 & 1 \end{pmatrix}.
\]
The lemma follows then from the well known fact that these two matrices generate $\SL(2,\F_\ell)$.
\end{proof}

\begin{proof}[Proof of the Theorem]
Lemma~\ref{lem:IV.3.2.1} shows that, for almost all $\ell$, $I_{v,\ell}$ and a fortiori $\widetilde{G}_\ell$ contains a transvection. On the other hand, we know (cf. \ref{sec:IV.2.1}) that $\widetilde{G}_\ell$ is irreducible for almost all $\ell$. Applying Lemma~\ref{lem:IV.3.2.2} to $\widetilde{G}_\ell$ we then see that $\widetilde{G}_\ell$ contains $\SL(E_\ell)$ for almost all $\ell$; hence we have (iv).
\end{proof}

\begin{remarks}
It seems likely that the condition ``$j$ is not integral'' can be replaced by the weaker one ``$E$ has no complex multiplication.'' \hspace{0.3em}$\rightarrow$ [yes: see \cite{serre-proprietes}.]
\end{remarks}

\subsechead{3.3}{Numerical example}

When $E$ is given explicitly and has a non-integral $j$, one may sometimes determine the finite set of $\ell$'s with $\widetilde{G}_\ell \neq \GL(2,\F_\ell)$. Take for instance $K=\Q$, and $E$ defined by the equation:
\[
  y^2 + x^3 + x^2 + x = 0.
\]
This is the curve $3^{+}$ of Ogg's list \cite{ogg}; its $j$ invariant is $2^{11}3^{-1}$, its discriminant is $\Delta = -2^4 \cdot 3$, its ``conductor'' is 24 (it is 2-isogenous to the modular curve $J_{24}$ corresponding to the congruence subgroup $\Gamma_0(24)$, cf. \cite{ogg}). The existence of a non-trivial 2-isogeny for $E$ shows that $\widetilde{G}_2 \neq \GL(2,\F_2)$ for $\ell=2$ ($\widetilde{G}_2$ is cyclic of order 2 and corresponds to the quadratic field $\Q(\sqrt{-3})$). \textbf{But, for $\ell \neq 2$, one has $\widetilde{G}_\ell = \GL(2,\F_\ell)$.} Indeed, $\widetilde{G}_\ell$ has the following properties:
\begin{enumerate}
\item[(a)] $\det(\widetilde{G}_\ell) = \F_\ell^{*}$, cf.~\ref{sec:IV.3.1}.
\item[(b)] $\widetilde{G}_\ell$ contains a transvection. This follows from Lemma~\ref{lem:IV.3.2.1} and the fact that $n$ is here equal to 1.
\item[(c)] $\widetilde{G}_\ell$ is irreducible. If not, there would be an isogeny $E \to E'$ of degree $\ell$ (defined over $\Q$). The curve $E'$ would have the same conductor 24 as $E$, hence would be one of the curves $1^{-},2^{+},3^{+},4^{-},5^{-},6^{+}$ of Ogg's list. But Ogg has proved that, for each such curve, there is an isogeny $E' \to E$ of degree 1, 2, 4 or 8. The map $E \to E' \to E$ would then be an endomorphism of $E$ of degree $\ell$, $2\ell$, $4\ell$ or $8\ell$, and this is impossible for $\ell \neq 2$ since $\End(E)=\Z$.
\end{enumerate}
Now, using lemma~\ref{lem:IV.3.2.2}, one sees that properties (a), (b), (c) imply that $\widetilde{G}_\ell = \GL(2,\F_\ell)$.

\begin{exercises}
Prove that $\widetilde{G}_\ell = \GL(2,\F_\ell)$ for all $\ell \neq 2$ when $K=\Q$ and $E$ is an elliptic curve of conductor $3 \cdot 2^{\lambda}$, where $\lambda \leq 6$. (Use Ogg's Table 1. For $\lambda=5$, note that the curves $7^{+}$ and $7^{-}$ become isomorphic over $\Q(i)$, but are not isogenous over $\Q$. For $\lambda=6$, use a similar argument, and observe that the curves $10^{+}$ and $18^{+}$ do not have the same number of points mod.\ 5, hence are not isogenous over $\Q$.) What happens when $\lambda=7,8$?
\end{exercises}

\subsechead{3.4}{Proof of the main lemma of \S\ref{sec:IV.3.1}}

We need first a few lemmas:

\begin{lemma}[1]\label{lem:IV.3.4.1}
Let $S_\ell = \mathrm{PSL}(2,\F_\ell) = \SL(2,\F_\ell)/\{\pm 1\}$, $\ell \geq 3$. Then $S_\ell$ is a simple group if $\ell \geq 5$. Every proper subgroup of $S_\ell$ is solvable or isomorphic to the alternating group $A_5$: the last possibility occurs only if $\ell \equiv \pm 1 \bmod 5$.
\end{lemma}

This is well known, cf. for instance Burnside \cite{burnside}, chap.~XX.

\begin{lemma}[2]\label{lem:IV.3.4.2}
No proper subgroup of $\SL(2,\F_\ell)$ maps onto $\mathrm{PSL}(2,\F_\ell)$.
\end{lemma}

This is clear for $\ell=2$, since $\mathrm{PSL}(2,\F_2) = \SL(2,\F_2)$. For $\ell \neq 2$, suppose there is such a proper subgroup $X$. We would then have $\SL(2,\F_\ell) = \{\pm 1\} \times X$, and this is absurd, since $\SL(2,\F_\ell)$ is generated by the elements $\begin{pmatrix} 1 & 1 \\ 0 & 1 \end{pmatrix}$ and $\begin{pmatrix} 1 & 0 \\ 1 & 1 \end{pmatrix}$ which are of order $\ell$, hence contained in $X$.

\begin{lemma}[3]\label{lem:IV.3.4.3}
Let $X$ be a closed subgroup of $\SL(2,\Z_\ell)$ whose image in $\SL(2,\F_\ell)$ is $\SL(2,\F_\ell)$. Assume $\ell \geq 5$. Then $X = \SL(2,\Z_\ell)$.
\end{lemma}

We prove by induction on $n$ that $X$ maps onto $\SL(2,\Z/\ell^n\Z)$. This is true for $n=1$. Assume it is true for $n$, and let us prove it for $n+1$. It is enough to show that, for any $s = \begin{pmatrix} a & b \\ c & d \end{pmatrix} \in \SL(2,\Z_\ell)$ which is congruent to 1 mod.\ $\ell^n$, there is $x \in X$ with $x \equiv s \bmod \ell^{n+1}$. Write $s = 1 + \ell^n u$; since $\det(s)=1$, one has $\operatorname{Tr}(u) \equiv 0 \bmod \ell$. But it is easy to see that any such $u$ is congruent mod.\ $\ell$ to a sum of matrices $u_i$ with $u_i^2=0$. Hence, we may assume that $u^2=0$. By the induction hypothesis, there exists $y \in X$ such that $y = 1 + \ell^{n-1}u + \ell^n v$, where $v$ has coefficients in $\Z_\ell$. Put $x = y^{\ell}$. We have:
\[
  x = 1 + \ell(\ell^{n-1}u + \ell^n v) + \binom{\ell}{2}(\ell^{n-1}u + \ell^n v)^2 + \dots + (\ell^{n-1}u + \ell^n v)^{\ell}.
\]
If $n \geq 2$, it is clear that $x \equiv 1 + \ell^n u \bmod \ell^{n+1}$. This is also true for $n=1$. Indeed, since $u^2=0$, and $u + \ell v \equiv u \bmod \ell$, we have
\[
  x \equiv 1 + \ell u + (u + \ell v)^{\ell} \bmod \ell^2.
\]
But $(u + \ell v)^2 \equiv \ell(u + vu) \bmod \ell^2$, hence:
\[
  (u + \ell v)^{\ell} \equiv \ell(u + vu)u^{\ell-2} \equiv 0 \bmod \ell^2
\]
since $\ell > 4$. This shows that $x \equiv 1 + \ell^n u \bmod \ell^{n+1}$ in all cases, and proves lemma~\ref{lem:IV.3.4.3}.

We now consider a closed subgroup $G$ of $X = \prod \GL(2,\Z_\ell)$ having the properties (a), (b), (c) of the main lemma of \ref{sec:IV.3.1}.

\begin{lemma}[4]\label{lem:IV.3.4.4}
Let $S$ be a finite set of primes, and $X_S = \prod_{\ell \in S} \GL(2,\Z_\ell)$. The image $G_S$ of $G$ by the projection $X \longrightarrow X_S$ is open in $X_S$.
\end{lemma}

Replacing $G$ by an open subgroup if necessary, we can assume that each $G_\ell$, $\ell \in S$, is contained in the group of elements congruent to 1 mod.\ $\ell$, hence that each $G_\ell$ is a pro-$\ell$-group. Since $G_S$ is a subgroup of $\prod_{\ell \in S} G_\ell$, it follows that $G_S$ is \textbf{pro-nilpotent} (projective limit of finite nilpotent groups), hence is the product of its Sylow subgroups. This shows that $G_S = \prod_{\ell \in S} G_\ell$, and since $G_\ell$ is open in $\GL(2,\Z_\ell)$ for $\ell \in S$, we see that $G_S$ is open in $X_S$.

Before we go further, we introduce some terminology. Let $Y$ be a profinite group, and $\Sigma$ a finite simple group. We say that $\Sigma$ \emph{occurs} in $Y$ if there exist closed subgroups $Y_1, Y_2$ of $Y$ such that $Y_1$ is normal in $Y_2$ and $Y_2/Y_1$ is isomorphic to $\Sigma$. We denote by $\operatorname{Occ}(Y)$ the set of classes of finite simple non-abelian groups occurring in $Y$. If $Y = \varprojlim Y_a$, and each $Y \to Y_a$ is surjective, we have
\[
  \operatorname{Occ}(Y) = \bigcup \operatorname{Occ}(Y_a).
\]
If $Y$ is an extension of $Y'$ and $Y''$, we have:
\[
  \operatorname{Occ}(Y) = \operatorname{Occ}(Y') \cup \operatorname{Occ}(Y'').
\]
Using these formulae and lemma~\ref{lem:IV.3.4.1}, one gets:
\[
  \operatorname{Occ}(\GL(2,\Z_\ell)) = \operatorname{Occ}(\SL(2,\Z_\ell)) = \operatorname{Occ}(S_\ell)
\]
where $S_\ell = \mathrm{PSL}(2,\F_\ell)$ as before, and:
\begin{align*}
  \operatorname{Occ}(S_\ell) &= \emptyset &&\text{if } \ell = 2,3, \\
  \operatorname{Occ}(S_\ell) &= \{S_\ell\} = \{A_5\} &&\text{if } \ell = 5, \\
  \operatorname{Occ}(S_\ell) &= \{S_\ell\} &&\text{if } \ell \equiv \pm 2 \bmod 5,\ \ell > 5, \\
  \operatorname{Occ}(S_\ell) &= \{S_\ell, A_5\} &&\text{if } \ell \equiv \pm 1 \bmod 5,\ \ell > 5.
\end{align*}

Let now $S$ be a finite set of primes so that $2,3,5 \in S$ and $\ell \notin S \Rightarrow \widetilde{G}_\ell \supset \SL(2,\F_\ell)$. Property (c) shows that such a set exists.

\begin{lemma}[5]\label{lem:IV.3.4.5}
The group $G$ contains $\prod_{\ell \notin S} \SL(2,\Z_\ell)$.
\end{lemma}

(This partial product is understood as a subgroup of the full product $X = \prod_{\ell} \GL(2,\Z_\ell)$.)

It is enough to show that $G$ contains each $\SL(2,\Z_\ell)$, $\ell \notin S$. Let $H_\ell = G \cap \GL(2,\Z_\ell)$. If $\ell \notin S$, the fact that $\widetilde{G}_\ell$ contains $\SL(2,\F_\ell)$ shows that $S_\ell \in \operatorname{Occ}(G_\ell)$ hence $S_\ell \in \operatorname{Occ}(G)$. On the other hand, $G/H_\ell$ is isomorphic to a closed subgroup of $\prod_{\ell} \GL(2,\Z_\ell)$, hence $S_\ell \notin \operatorname{Occ}(G/H_\ell)$ (we use the obvious fact that the simple groups $S_p$, $p \geq 5$, are pairwise non-isomorphic). Since $\operatorname{Occ}(G) = \operatorname{Occ}(H_\ell) \cup \operatorname{Occ}(G/H_\ell)$, we then have $S_\ell \in \operatorname{Occ}(H_\ell)$. Let $\widetilde{H}_\ell$ be the image of $H_\ell$ in $\SL(2,\F_\ell)$; the kernel of $H_\ell \to \widetilde{H}_\ell$ being a pro-$\ell$-group, we have $\operatorname{Occ}(H_\ell) = \operatorname{Occ}(\widetilde{H}_\ell)$, hence $S_\ell \in \operatorname{Occ}(\widetilde{H}_\ell)$. Hence $\widetilde{H}_\ell$ maps onto $S_\ell = \mathrm{PSL}(2,\F_\ell)$, and, by lemma~\ref{lem:IV.3.4.2}, we have $\widetilde{H}_\ell = \SL(2,\F_\ell)$ and, by lemma~\ref{lem:IV.3.4.3}, $H_\ell = \SL(2,\Z_\ell)$. Hence $G$ contains $\SL(2,\Z_\ell)$.

\begin{lemma}[6]\label{lem:IV.3.4.6}
The group $G$ contains an open subgroup of $\prod_{\ell} \SL(2,\Z_\ell)$.
\end{lemma}

Let $S$ be as in lemma~\ref{lem:IV.3.4.5}; let $G_S$ be the projection of $G$ into $\prod_{\ell \in S} \GL(2,\Z_\ell)$ and $G'_S$ the projection into the complementary product $\prod_{\ell \notin S} \GL(2,\Z_\ell)$. Let $H_S = G \cap \prod_{\ell \in S} \GL(2,\Z_\ell)$ and $H'_S = G \cap \prod_{\ell \notin S} \GL(2,\Z_\ell)$, so that $H_S \subset G_S$, $H'_S \subset G'_S$. One has canonical isomorphisms:
\[
  G_S / H_S \simeq G / (H_S \times H'_S) \simeq G'_S / H'_S.
\]
Lemma~\ref{lem:IV.3.4.5} shows that $H'_S$ contains $\prod_{\ell \notin S} \SL(2,\Z_\ell)$, so that $G'_S / H'_S$ is abelian. Hence $G_S / H_S$ is abelian and $H_S$ contains the closure $(G_S,G_S)$ of the commutator group of $G_S$. By lemma~\ref{lem:IV.3.4.4}, $G_S$ is open in $\prod_{\ell \in S} \GL(2,\Z_\ell)$. It is easy to see that this implies that $(G_S,G_S)$ contains an open subgroup of $\prod_{\ell \in S} \SL(2,\Z_\ell)$ (this follows for instance from the fact that the derived Lie algebra of $\mathfrak{gl}_2$ is $\mathfrak{sl}_2$). Hence $H_S$ contains an open subgroup $U$ of $\prod_{\ell \in S} \SL(2,\Z_\ell)$. Using lemma~\ref{lem:IV.3.4.5}, we then see that $G$ contains $U \times \prod_{\ell \notin S} \SL(2,\Z_\ell)$ which is open in $\prod_{\ell} \SL(2,\Z_\ell)$.

\begin{proof}[End of the proof]
Consider the determinant map
\[
  \det : \prod_{\ell} \GL(2,\Z_\ell) \longrightarrow \prod_{\ell} \Z_\ell^{*},
\]
whose kernel is $\prod \SL(2,\Z_\ell)$. Hypothesis (c) means that the image of $G$ by this map is open and lemma~\ref{lem:IV.3.4.6} shows that $G \cap \operatorname{Ker}(\det)$ is open in $\operatorname{Ker}(\det)$. This is enough to imply that $G$ itself is open, q.e.d.
\end{proof}

\begin{exercises}
\begin{enumerate}
\item $\mathrm{a)}$ Generalize lemma~\ref{lem:IV.3.4.3} to $\SL(d,\Z_\ell)$ for $d \geq 2$, $\ell \geq 5$ (same method).
\item $\mathrm{b)}$ Show that the only closed subgroup of $\SL(d,\Z_3)$ which maps onto $\SL(d,\Z/3^2\Z)$ is $\SL(d,\Z_3)$ itself.
\item $\mathrm{c)}$ Show that the only closed subgroup of $\SL(d,\Z_2)$ which maps onto $\SL(d,\Z/2^3\Z)$ is $\SL(d,\Z_2)$ itself.
\item Let $E$ be the unramified quadratic extension of $\Q_2$, and $\mathcal{O}_E$ its ring of integers. Let $x \mapsto \overline{x}$ be the non-trivial automorphism of $E$.
\item $\mathrm{a)}$ Show that $\mathcal{O}_E$ contains a primitive third root of unity $z$.
\item $\mathrm{b)}$ Show that $\mathcal{O}_E$ contains an element $u$ with $u \cdot \overline{u} = -1$ (take for instance $u = (1+\sqrt{5})/2$).
\item $\mathrm{c)}$ Let $\alpha$ and $\beta$ be the $\Z_2$-linear endomorphisms defined by $\alpha(x) = zx$, $\beta(x) = \overline{u}x$, where $z$ and $u$ are as in a), b) above. Show that $\alpha$ is of order 3, $\beta$ of order 4, and $\beta \alpha \beta^{-1} = \alpha^{-1}$, so that $\alpha$ and $\beta$ generate a non-abelian group $G$ of order 12.
\item $\mathrm{d)}$ Show that $G$ is contained in $\SL(\mathcal{O}_E) \simeq \SL(2,\Z_2)$ and that reduction mod.\ 2 defines a homomorphism of $G$ onto $\SL(2,\F_2)$. (Hence lemma~\ref{lem:IV.3.4.3} does not extend to the case $\ell = 2$.)
\item Let $S_9 = \SL(2,\Z/9\Z)$, $S_3 = \SL(2,\Z/3\Z)$ and $g = \operatorname{Ker}(S_9 \longrightarrow S_3)$. The group $g$ is isomorphic to a three-dimensional vector space over $\F_3$. Let $x \in \mathrm{H}^2(S_3, g)$ be the cohomology class corresponding to the extension
\[
  1 \longrightarrow g \longrightarrow S_9 \longrightarrow S_3 \longrightarrow 1.
\]
\item $\mathrm{a)}$ Show that the restriction of $x$ to a 3-Sylow subgroup of $S_3$ is zero (note that $\SL(2,\Z)$ contains an element of order 3, viz.\ $\begin{pmatrix} 0 & -1 \\ 1 & -2 \end{pmatrix}$).
\item $\mathrm{b)}$ Deduce from a) that $x = 0$, i.e.\ that there exists a subgroup $X$ of $S_9$ which is mapped isomorphically onto $S_3$. (The inverse image of $X$ in $\SL(2,\Z_3)$ is a non-trivial subgroup which is mapped onto $S_3$; hence lemma~\ref{lem:IV.3.4.3} does not extend to the case $\ell = 3$.)
\end{enumerate}
\end{exercises}
